% Generated draft: requires mathematical and source review.
\documentclass[11pt]{article}
\usepackage[margin=25mm]{geometry}
\usepackage{fontspec}
\setmainfont{Latin Modern Roman}
\usepackage{amsmath,amssymb,mathtools,unicode-math}
\setmathfont{Latin Modern Math}
\usepackage{enumitem,needspace,xurl,hyperref}
\setlength{\parindent}{0pt}
\setlength{\parskip}{5pt}
\setlength{\emergencystretch}{3em}
\newcommand{\N}{\mathbb{N}}
\newcommand{\Z}{\mathbb{Z}}
\newcommand{\Q}{\mathbb{Q}}
\newcommand{\R}{\mathbb{R}}
\newcommand{\C}{\mathbb{C}}
\newcommand{\F}{\mathbb{F}}
\newcommand{\A}{\mathbb{A}}
\newcommand{\PP}{\mathbb{P}}
\newcommand{\E}{\mathbb{E}}
\newcommand{\eps}{\varepsilon}
\newcommand{\dd}{\,\mathrm{d}}
\begin{document}
\section*{UnsolvedMath ID: 20001292}
\textbf{Source problem code:} catalogue entry
\subsection{Definitions and mathematical statement}
David Jerison's source asks whether a least-area interface dividing a symmetric convex body into halves, or any fixed proportion, is a graph, although it need not be a hyperplane. A precise relative-isoperimetric interpretation takes a centrally symmetric convex body $K\subset\R^n$, $0<\theta<1$, and a finite-perimeter set $E\subset K$ minimizing
\[
P(E;\operatorname{int}K)=|D\mathbf1_E|(\operatorname{int}K),\qquad |E|=\theta|K|.
\]
Only interior interface area is counted, not the part on $\partial K$. The graph question concerns the interior measure-theoretic interface, after some rotation and with an appropriate representative. The original does not fix regularity, treatment of singular points or the exact graph convention.

The attached aid settles a planar half-volume family. Let $R=(-a,a)\times(-b,b)$ with $a\ge b>0$. For every finite-perimeter $E\subset R$ with $|E|=2ab$,
\[
P(E;R)\ge2b.
\]
If $a>b$, equality holds exactly, up to planar null sets, for the left or right half-rectangle $R\cap\{x<0\}$ and $R\cap\{x>0\}$. If $a=b$, the top and bottom halves are the only additional equality cases. Their reduced interior interfaces are the corresponding central shortest-length segments and global affine graphs over intervals after rotation. This does not establish the arbitrary-proportion or general-convex-body assertion.

The aid describes a BV-compatible even reflection into the $4a$ by $4b$ flat torus, multiplying area and relative perimeter by four without seam jumps, followed by a classification of half-area equality bands and imposition of the fold symmetries. Its full torus equality proof is not supplied here.

The original remark discusses the Simons cone with displayed equation
\[
x_1^2+x_2^2+x_3^2+x_4^2=x_5^2+x_6^2+x_7^2+x_8^2.
\]
It reports minimization with a prescribed boundary on the ball, nongraph behavior, and smaller area for a bisecting plane. The fixed-boundary comparison differs from unconstrained relative half-volume minimization, so it is not a counterexample to the question. The remark further claims a growing-dimension approach to the plane area and asks whether a Gaussian limit could conflict with Borell's finite-dimensional uniqueness of half-spaces. The cone family, normalization and limiting notion are not specified in this record; that proposed limit and its inference are retained as an unverified source discussion, not an established contradiction or failure of Borell's theorem.
\subsection{Short English statement}
Ask whether relative area-minimizing partitions of symmetric convex bodies have graphical interfaces. The supplied aid proves the exact half-area rectangle case, while preserving the different fixed-boundary and Gaussian-limit qualifications in the original Simons-cone remark.
\subsection{Sources}
\begin{itemize}
\item Local UnsolvedMath registry, permanent ID $20001292$, source code catalogue entry; canonical \texttt{aim-convex-geometry-notes.json}, \texttt{notes[23]}.
\item AIM \emph{Problems from the Workshop on Algorithmic Convex Geometry}, version dated 31 October 2007, problem 5 attributed to David Jerison: \url{https://aimath.org/WWN/convexgeometry/convexgeometry.pdf}. Its remark names Simons and Borell without exact references. No PDF inspection is claimed here.
\item Attached research attempt 1: partial result, recorded run status open, rigor proved, novelty confidence low; registry status partially solved, classification PARTIAL-PROGRESS. Its rectangle theorem and reflection argument are machine generated and not peer reviewed; independently verify before citation.
\end{itemize}
\end{document}
